\begin{abstract}
\iffalse

模糊測試（fuzzing）作為資訊安全領域的重要技術，主要透過提供大量隨機輸入，以找出目標程式中的潛在漏洞與錯誤。然而，現有的模糊測試工具多數只針對固定的命令列選項進行輸入檔或標準輸入的變異，對於多選項（多參數）的探索缺乏靈活性且需要專家介入定義參數組態檔。為解決這一限制，本研究提出了一個名為 KOFTA 的選項感知模糊測試架構。KOFTA 結合靜態分析與動態分析自動化地探索可用的命令列選項與選項參數，並利用污染源推論技術精確找出影響分支路徑的關鍵參數，以充分探索廣闊的輸入空間。此外，透過改良 AFL 的 forkserver，使其能夠動態更換目標程式的命令列參數，有效降低測試過程中的系統呼叫成本。實驗結果顯示，相較於傳統的多參數模糊測試工具，KOFTA 在提升測試覆蓋率、發掘未知選項與參數，以及漏洞挖掘方面展現了明顯的優勢，並提供了一個開源的原型以供未來研究與應用。
\else

Fuzzing, a crucial technique in information security, primarily identifies potential vulnerabilities and errors in target programs by providing a large amount of random inputs. However, most existing fuzzing tools only mutate input files or standard inputs for fixed command-line options, lacking flexibility in exploring multiple options (multiple arguments) and requiring expert intervention to define option configuration files. To address this limitation, this research proposes an option-aware fuzzing framework called KOFTA. KOFTA combines static and dynamic analysis to automatically explore available command-line options and parameters, and utilizes taint inference technique to accurately identify critical parameters affecting branch paths, thus fully exploring a broad input space. Furthermore, by improving AFL's forkserver, it enables dynamic replacement of the target program's command-line arguments, effectively reducing system call costs during testing. Experimental results show that compared to traditional multi-argument fuzzing tools, KOFTA demonstrates significant advantages in improving test coverage, discovering undocumented options and parameters, and vulnerability detection. An open-source prototype is also provided for future research and application.

\fi
  \begin{keywords}
  Fuzzing, Multi-Option Fuzzing, Option-Aware Fuzzing, Taint Analysis, Taint Inference
  \end{keywords}
\end{abstract}